\section{The \texttt{MM\_h} McXtrace Component}
Date: Feb. 2017
Version: 1.1
Release: McXtrace 1.2
Origin: DTU Physics, DTU Space

Single Pore as part of the Silicon Pore Optics (SPO) as envisioned for the ATHENA+ space telescope.

\subsection*{Identification}
\begin{itemize}
  \item \textbf{Author:} Erik B Knudsen and Desiree D. M. Ferreira
  \item \textbf{Origin:} DTU Physics, DTU Space
  \item \textbf{Date:} Feb. 2016
\end{itemize}

\subsection*{Description}
A single pore is simulated, which may have thick walls. The top and bottom are curved cylindrically azimuthally, and according to the Wolter I optic lengthwise (sagitally). This is the hyperbolic part. The azimuthal curvature is defined by the radius parameters.

To intersect the Wolter I plates we take advatage of the azimuthal symmetry and only consider the radial component of the photon's wavevector.

Example: MM\_h( pore\_th=0, ring\_nr=1, Z0=12, pore\_width=0.83e-3, pore\_height=0.605e-3, chamfer\_width=0.17e-3, mirror\_reflec="mirror\_coating\_unity.txt", R\_d=0, size\_file="ref\_design\_breaks.txt")

\subsection*{Input parameters}
Parameters in \textbf{boldface} are required; the others are optional.

\begin{longtable}{p{0.22\textwidth}p{0.12\textwidth}p{0.46\textwidth}p{0.14\textwidth}}
\toprule
\textbf{Name} & \textbf{Unit} & \textbf{Description} & \textbf{Default} \\
\midrule
\endhead
\textbf{pore\_th} &  &  &  \\
\textbf{ring\_nr} &  &  &  \\
\textbf{Z0} & m & Distance between intersection plane and the focal spot (essentially focal length). &  \\
\textbf{pore\_height} & m & Height of the pore. (Hence the inner radii are radius\_\{m,h\}-pore\_height) &  \\
\textbf{pore\_width} &  &  &  \\
chamfer\_width &  &  & 0 \\
gap & m & Gap between intersection with parabolic section and actual plate. & 0 \\
zdepth &  &  & 0 \\
mirror\_reflec &   & Data file containing reflectivities of the reflector surface (TOP). & "" \\
bottom\_reflec &   & Data file containing reflectivities of the bottom surface (BOTTOM). & "" \\
side\_reflec &   & Data file containing reflectivities of the side walls (LEFT and RIGHT). & "" \\
size\_file &  &  & "" \\
non\_specular\_file &  &  & "" \\
R\_d &   & Default reflectivity value to use if no reflectivity file is given. Useful f.i. is one surface is reflecting and the others absorbing. & 1 \\
waviness &  &  & 0 \\
longw &  &  & 0 \\
\bottomrule
\end{longtable}

\subsection*{Links}
\begin{itemize}
  \item Component source code found in file \texttt{MM\_h.comp}.
\end{itemize}
\IfFileExists{astrox/MM_h_static.tex}{\input{astrox/MM_h_static.tex}}{}